\chapter{Counting Is Barely More Than Telling Apart}

The instrument can tell two things apart. That was the whole of the last chapter, and it seemed like almost
nothing --- a bare difference, a this that is not a that, with no quantity anywhere in it. Now watch how little it
takes to get from there to number. Not a leap, not a new faculty granted from outside, not some fresh power the
machine did not have a moment ago. Just one small habit added to the one it already has: keep telling things
apart, and keep track of how many times you have had to. That is counting. It is the second thing the instrument
can do, and it is built directly on the first, and it is barely more than the first. A number, in this machine,
is not a thing that arrives with its own nature and its own laws; it is a tally of differences admitted, and it
means nothing more than that.

We are still, be warned, at the bottom of the tower. Nothing here is measured, nothing is weighed, nothing has a
unit. The count the instrument keeps is not yet a length or a duration or a charge; it is a pure how-many, the
plainest quantity there is, and it will be many chapters before it wears any physics. But the plainest quantity
is where number has to begin if it is to begin honestly, from below, without being handed down whole. So we take
the second step exactly as we took the first: we add only what the construction earns, and we watch precisely
where the count comes from, because where it comes from is the whole of what it will ever mean.

There is a second reason to go slowly here, beyond honesty. The step from difference to number is the one step in
the whole construction that everyone thinks they already understand, and precisely because it feels obvious it is
the easiest place to smuggle something in. We are so used to numbers arriving whole --- with their line, their
arithmetic, their sense of size all attached --- that it takes an effort to see how little of that the machine
actually has at this stage, and how much of it will have to be built, later, rung by rung, at a cost we will
insist on watching. So the chapter is deliberately meager. It gives the instrument the tally and the order and
not one thing more, and it asks you to notice the meagerness, because the meagerness is the guarantee: what is
not yet in the number cannot later be claimed from it for free.

\section{Number from a Difference Admitted}

The claim is that counting is one short step past telling apart: the instrument builds a count by keeping a
running tally of the differences it admits.

Against the machine, it comes to this. On top of the first rung --- the one that tells two marks apart --- the
construction sets a second rung, and a third close behind it. The second gives the machine a process for taking a
new mark and \emph{admitting} it: deciding whether this difference is one to be entered into the account at all,
and, if it is, advancing the tally by one. The third gives the machine an ordering --- a way of holding its
tallies in a line, so that one count stands before another, smaller before larger. Put the two together and you
have the whole apparatus of counting: a way to say \emph{admit this} and a way to say \emph{that makes one more
than before}. There is nothing else in it. The machine does not reach for the numbers as a pre-existing supply
and pull one down; it grows each count from the last by the single act of admitting one more difference.

Follow the number itself for a moment, because it is built as plainly as everything else here. It starts at
nothing --- no differences admitted, no tally. Admit one, and the tally becomes one. Admit another, and it
becomes one more than that. Each number is literally the one before it with a single admission laid on top, and
that is all a number is: a stack of admissions, each resting on the one beneath. When the machine wants a larger
count it does not compute it or look it up; it climbs, one admitted difference at a time, from nothing. The whole
of arithmetic that will eventually stand on this begins as a tower of single steps, each step the same small act
the first rung already performed --- noticing that here is one more thing, not the same thing again.

It is worth pausing on how strange, and how strong, this plainness is. The instrument does not know what a number
\emph{is} in any sense beyond this. It has no notion of number as a point on a line, no notion of number as a
size or a magnitude or a position; it has only the tally and the act of adding to it. And yet everything the word
number will ever mean, in this book, is going to be built out of exactly this and nothing richer. When, chapters
from now, a count becomes a charge, or a mass, or the reading on a gauge, it will not have acquired some new
substance along the way. It will still be a tally of differences admitted --- the same stack of single steps ---
wearing a physical name. That is the discipline paying off in advance: because we refused to let number arrive
with more than the tally in it, we will be able to say later, of every quantity the instrument reports, exactly
what it is made of and exactly what it cost. A number that came pre-loaded with meaning could never be audited
that way. This one can, because there is nothing in it we did not watch go in.

And notice, as before, what is \emph{not} here. The machine cannot yet subtract, cannot yet compare two counts
for how far apart they are, cannot yet do anything but build a tally up and hold tallies in order. Those powers
come later, each on its own rung, each earned. What it has now is exactly enough to say how many differences it
has admitted, and to say which of two such how-manys is the greater. That is counting, and it is the second
thing, and it was almost free.

\section{The Sign on the Order}

There is one feature of the ordering that deserves to be pulled out and looked at, quietly, because it is doing
more than it lets on. When the machine holds two tallies and asks which stands before the other, it does not ask
a bare arithmetical question. It asks the question with a \emph{sign convention} attached --- a rule about which
direction counts as growth.

Read against the machine, the ordering is built in two branches, not one. When two readings agree in their
underlying truth --- when the difference being counted is, in the relevant sense, the same kind of difference ---
the order runs the ordinary way: more admissions means a larger count, growth is upward, and the tally that has
admitted more stands later in the line. But the machine also carries a second branch, for the case where the two
readings \emph{disagree} in their underlying truth, and there the order runs the other way: what counts as growth
is reversed, and the tally with fewer admissions is the one held to be larger. The machine decides which branch
it is in by consulting the truth beneath the count --- the same decidable truth from the first rung --- and
orders accordingly.

This looks, at first, like a fussy technicality, a bookkeeper's footnote about which way is up. It is not. It is
the first appearance, in the plainest possible setting, of a distinction that will run through the entire book
and eventually carry a great deal of physics on its back: the distinction between a count that runs \emph{with}
the thing being counted and a count that runs \emph{against} it. Here it is only a sign on an ordering ---
whether admitting a difference makes the number go up or down, depending on whether the difference agrees with
the account or cuts against it. Later it will be the difference between two whole ways of reading a record, and
the machine will have a name for each way. We are nowhere near that yet, and it would be dishonest to pretend the
name is earned here. But the seed is here, in the branch on the order, and it is worth marking the spot so that
when the two ways of reading arrive in full you will recognize where they were first quietly planted: in the
question of which direction, on a plain tally of differences, counts as more.

To feel the two branches without any physics on them, take the barest illustration. Suppose the machine is
tallying differences that all agree in their underlying truth --- all pointing the same way, all of a kind. Then
each new admission simply pushes the count higher, and to have admitted more is straightforwardly to be larger;
the order is the familiar one, and there is nothing to remark. But now suppose a difference comes in that cuts
against the grain of the account --- one whose underlying truth runs opposite to the rest. The machine still
admits it, still advances a tally; but the branch it lands in reverses the reckoning, so that this admission,
counted honestly, tells \emph{against} the total rather than for it. Nothing mysterious has happened. The machine
has only kept faith with the truth beneath each count, and that truth has two possible orientations, and the
order respects both. The point to carry forward is only this: from the very first tally, growth already has a
direction that depends on what is being counted, and the machine reads that direction off the world rather than
imposing one. The two ways of counting are already latent in the plainest count there is.

\section{The Floor, and What the Count Is Of}

The floor first, because the discipline holds and the honesty of the whole book is in holding it. Nothing in this
chapter has been proved. Counting, and the ordering with its sign, are \emph{built} --- laid down as the machine's
construction, not established by an argument that could have failed. There was no lurking claim here that might
have gone the other way; there was only a decision to equip the machine with a tally and an order. The number
descends from the first difference by definition, and the grade is the same as the last chapter's: what is
established, is established by construction. When the first genuine proof arrives --- and it is close now, closer
than you might think --- you will feel the ground change. It has not changed yet.

Now the reading, which is short because the chapter was short. A count, in this machine, is a how-many of
differences admitted, and that fixes what a count is \emph{of}. It is of the world, in exactly the sense the last
chapter drew: the differences the machine admits are differences that are really there to be told, and the tally
of them is not the machine's invention but its record. That there are three of something rather than two is not a
fact the instrument authored; it counted, and the world supplied the differences to count. And yet the account of
the count --- the direction the order runs, the sign on the growth, the choice of which branch to be in --- has a
part that is the machine's own. The how-many is the world's; the which-way-is-up is, in part, ours. You can hear
the book's one line again, one rung higher than we left it: the count beneath is the world's, and the convention
we wrap it in --- the sign, the direction, the bookkeeping of growth --- is the hand's. Counting was barely more
than telling apart. It kept, as it had to, the same honest seam down the middle: the difference is given, and the
way we choose to face it is ours.
